\documentclass[10pt,a4paper]{amsart}
\usepackage[margin=19mm]{geometry}
\usepackage{amsmath,amssymb,mathtools}
\usepackage[scr=rsfs,cal=euler]{mathalfa}
\usepackage{xfrac}
\usepackage[unicode,hidelinks,bookmarks=false]{hyperref}
\newcommand{\ie}{i.e.\ }
\newcommand{\cf}{cf.\ }
\newcommand{\resp}{respectively\ }
\newcommand{\Subsection}{\subsection}
\providecommand{\nobreakdash}{\nobreak\hbox{-}}
\setlength{\emergencystretch}{2em}
\multlinegap0pt
\usepackage{amssymb}
\newtheorem{lemma}{Lemma}
\title{How large part of a graph can be covered\\ by the neighborhoods of $k$ vertices?}
\author{Janos Pach}
\date{Source: arXiv:2604.27639v1 (CC BY 4.0)}
\begin{document}
\maketitle

\noindent\emph{Source and licence.} How large part of a graph can be covered\\ by the neighborhoods of $k$ vertices?. Version arXiv:2604.27639v1 (CC BY 4.0), distributed by arXiv under the Creative Commons Attribution 4.0 International licence, http://creativecommons.org/licenses/by/4.0/, as recorded in the arXiv licence metadata for this exact version and shown on the abstract page https://arxiv.org/abs/2604.27639v1. Full licence notice: \texttt{licenses/arxiv-2604.27639-CC-BY-4.0.txt}.\par\smallskip

\noindent\textit{Editorial scope.} Counted unit: the article's lemma lemma2 (source0.tex lines 206--211). The article's complete original proof of it is delivered with it (source0.tex lines 213--349, one \texttt{proof} environment of 137 lines). No mathematics is written, repaired or re-derived: the statement and the proof are the article's own lines, in the article's own notation, and no retained line is substituted or re-flowed.  No further source-local context is needed: every cross-reference of every retained line resolves inside the two delivered spans.  Nothing was omitted from any retained span, so the package carries no omission record.  No retained line contains a citation, so the wrapper carries no bibliography and no bibliography entry.  No retained line contains a figure environment, an image-inclusion command or drawing code, so no image is delivered and the package declares no asset dependency.  Editorial formatting only, in addition: an AMS \texttt{amsart} preamble that restates the article's own declarations and macros used by the retained lines, namely the environments \texttt{lemma} on separate counters, together with the standard packages those lines need.  The retained environments print in this document as Lemma 1 in order of appearance, whereas the article prints the same items with the numbering of its own full document.  The complete native source of the article as distributed in the arXiv e-print for this exact version is bundled unchanged as \texttt{sources/199488/source0.tex}.
\begin{lemma}\label{lemma2}
Let $k\ge 2$ be fixed. For every $\;0<c<1$ and every $s\in[c,\sqrt c]$, we have
\[
F(s)\ge \min\{G(c),\sqrt c\}.
\]
\end{lemma}

\begin{proof}
We split the proof into two cases.

\medskip
\noindent\textbf{Case 1: $G(c)\le \sqrt c$.}
\smallskip

It will be convenient to set a new variable
\[
x:=1-s,
\]
so that we have $x\in[1-\sqrt c,\,1-c]$ and
\[
F(s)=(1-x)x^{k-1}+c\,\frac{1-x^{k-1}}{1-x}=(1-x)x^{k-1}+c\sum_{i=0}^{k-2}x^i.
\]
Hence,
\begin{align*}
F(s)-G(c)
&=x^{k-1}-x^k+\sum_{i=0}^{k-2}x^i-(1-c)\sum_{i=0}^{k-2}x^i-1+(1-c)^k \\
&=(1-c)^k-(1-c)\sum_{i=0}^{k-2}x^i+\sum_{i=1}^{k-1}x^i-x^k \\
&=\bigl((1-c)^k-x^k\bigr)+\bigl(x-(1-c)\bigr)\sum_{i=0}^{k-2}x^i \\
&=\bigl((1-c)-x\bigr)\left(\sum_{i=0}^{k-1}(1-c)^{k-1-i}x^i-\sum_{i=0}^{k-2}x^i\right).
\end{align*}

We need to show that $F(s)-G(c)\ge 0$. Observe that $1-c-x\ge0.$ If $1-c=x,$ then we are done. Otherwise,
\[
\frac{F(s)-G(c)}{1-c-x}=\sum_{i=0}^{k-1}(1-c)^{k-1-i}x^i-\sum_{i=0}^{k-2}x^i.
\]
The right-hand side is a polynomial $P(x)$ of degree $k-1$ in $x$, with coefficients
\[
(1-c)^{k-1}-1,\ (1-c)^{k-2}-1,\ \dots,\ (1-c)-1,\ 1.
\]
Since $0<1-c<1$, this sequence has exactly one sign change. By Descartes' rule of signs, $P(x)$ has at most one positive root.
In view of the fact that $P(0)=(1-c)^{k-1}-1<0,$ it follows that if a positive root exists, then $P(x)<0$ before the root and $P(x)>0$ after that.


Evaluating $P$ at $x_0:=1-\sqrt c$ and taking into account our assumption $G(c)\le \sqrt c$,  we find that
\[
P(x_0)=\frac{F(\sqrt c)-G(c)}{1-c-x_0}=\sqrt{c}-G(c)\ge 0.
\]
Hence, $P(x)\ge 0$ for every $x\in[x_0,1-c]$, which implies that
$F(s)-G(c)\ge 0,$ as required.

\medskip
\noindent\textbf{Case 2: $G(c)\ge \sqrt c$.}
\smallskip

As before, we set $x:=1-s,$ so that $x\in[1-\sqrt{c},\,1-c]$. Now we need to show that $F(s)\ge\sqrt c$.

%So
%\[
%F(s)-\sqrt{c}=sx^{k-1}+c\sum_{i=0}^{k-2}x^i-\sqrt{c}.
%\]
%Now define
%\[
%Q(x):=\sqrt{c}\sum_{i=0}^{k-2}x^i-x^{k-1}.
%\]
%We claim that
%\[
%F(s)-\sqrt{c}=(\sqrt{c}-s)Q(x).
%\]
%Indeed,
%\begin{align*}
%(\sqrt{c}-s)Q(x)
%&=(\sqrt{c}-s)\left(\sqrt{c}\sum_{i=0}^{k-2}x^i-x^{k-1}\right) \\
%&=\sqrt{c}(\sqrt{c}-s)\sum_{i=0}^{k-2}x^i-(\sqrt{c}-s)x^{k-1} \\
%&=c\sum_{i=0}^{k-2}x^i-\sqrt{c}s\sum_{i=0}^{k-2}x^i-\sqrt{c}x^{k-1}+sx^{k-1}.
%\end{align*}
%Now use the geometric-sum identity
%\[
%s\sum_{i=0}^{k-2}x^i=(1-x)\sum_{i=0}^{k-2}x^i=1-x^{k-1}.
%\]
%Therefore
%\[
%\sqrt{c}s\sum_{i=0}^{k-2}x^i=\sqrt{c}(1-x^{k-1}).
%\]
%Substituting this above gives
%\begin{align*}
%(\sqrt{c}-s)Q(x)
%&=c\sum_{i=0}^{k-2}x^i-\sqrt{c}(1-x^{k-1})-\sqrt{c}x^{k-1}+sx^{k-1} \\
%&=c\sum_{i=0}^{k-2}x^i-\sqrt{c}+sx^{k-1} \\
%&=F(s)-\sqrt{c}.
%\end{align*}

A direct calculation, similar to the one used in Case 1, gives that
\begin{equation}\label{eq3}
F(s)-\sqrt c=(\sqrt c-s)Q(x),
\end{equation}
where
\[
Q(x):=\sqrt{c}\sum_{i=0}^{k-2}x^i-x^{k-1}
\]
is a polynomial of degree $k-1$ is $x$.

Since $\sqrt{c}-s\ge 0$, it remains to show that $Q(x)\ge 0$ on the interval $[1-\sqrt{c},1-c]$.

The coefficients of $Q(x)$ are
\[
\sqrt{c},\ \sqrt{c},\ \dots,\ \sqrt{c},\ -1,
\]
so there is exactly one sign change.
By Descartes' rule of signs, $Q(x)$ has at most one positive root.
We have
$Q(0)=\sqrt{c}>0,$
and the leading coefficient is $-1$. Thus, if a positive root exists, then $Q(x)>0$ before the root and $Q(x)<0$ after that.
Now let $x_0:=1-c.$
We compute
\begin{align*}
Q(x_0)
&=\sqrt{c}\sum_{i=0}^{k-2}(1-c)^i-(1-c)^{k-1}=\sqrt{c}\cdot \frac{1-(1-c)^{k-1}}{c}-(1-c)^{k-1} \\
&=\frac{1-(1-c)^{k-1}}{\sqrt{c}}-(1-c)^{k-1}=\frac{1-(1-c)^{k-1}(1+\sqrt{c})}{\sqrt{c}}.
\end{align*}
On the other hand,
\begin{align*}
G(c)\ge \sqrt{c}
&\iff 1-(1-c)^k\ge \sqrt{c} \\
%&\iff (1-c)^k\le 1-\sqrt{c} \\
%&\iff (1-c)^{k-1}(1-c)\le 1-\sqrt{c} \\
&\iff (1-c)^{k-1}(1-\sqrt{c})(1+\sqrt{c})\le 1-\sqrt{c} \\
&\iff (1-c)^{k-1}(1+\sqrt{c})\le 1,
\end{align*}
Thus, our assumption $G(c)\ge \sqrt{c}$ is equivalent to the statement that
$Q(x_0)\ge 0.$
Therefore, any positive root of $Q$ is at least $x_0$, so $Q(x)\ge 0$ for all $x\in[0,x_0]$.
In particular, we have $Q(x)\ge 0$ for all $x\in[1-\sqrt{c},1-c]$. In view of (\ref{eq3}), this implies that
\[
F(s)-\sqrt{c}\ge 0.
\]
as required.
\medskip

Combining the above two cases yields
\[
F(s)\ge \min\{G(c),\sqrt c\},
\]
which completes the proof of Lemma~\ref{lemma2} and, hence, the theorem.
\end{proof}

\end{document}
